% 07-recomendaciones.tex — Sección 7: recomendaciones técnicas, derivadas de los
% datos de este experimento, no de opiniones generales sobre los motores.

\section{Recomendaciones técnicas para despliegues documentales en la nube}
\label{sec:recomendaciones}

Las recomendaciones salen de lo medido aquí, con sus límites: un solo conjunto de
datos de \num{500000} documentos, dos consultas y un host virtualizado. No son
afirmaciones generales sobre los dos motores.

\paragraph{1. Elegir el motor por la forma de la consulta, no por el motor.}
Es el resultado más claro: MongoDB gana \texttt{Q1} en las nueve celdas y
PostgreSQL gana \texttt{Q2} en ocho de nueve. Si la carga dominante desanida
arreglos y agrupa, MongoDB rinde de 2 a 5 veces más; si ordena y se queda con un
Top-K, PostgreSQL rinde de 1.2 a 10.7 veces más. Un despliegue que haga las dos
cosas no se resuelve eligiendo motor, sino decidiendo qué consulta importa más.

\paragraph{2. Dimensionar \texttt{work\_mem} mirando el plan, no el tamaño de la tabla.}
\texttt{Q2} ordena \num{175049} filas y necesita \textbf{273~kB}, porque el
\texttt{LIMIT} la convierte en un \emph{top-K}. \texttt{Q1} sobre los mismos datos
necesita unos 0.8~GB \textbf{por consulta} y vuelca a disco con cualquier valor
razonable de \texttt{work\_mem}. Subir el parámetro de 4 a 64~MB no cambió nada
medible en \texttt{Q2} y no habría salvado a \texttt{Q1}: lo que decide es la forma
del plan.

\paragraph{3. En PostgreSQL, vigilar \texttt{temp\_bytes} como métrica de primera línea.}
El volcado no se nota en la CPU ni en el RSS: se nota en \textbf{\cifraTempTotalTB{}~TB}
escritos en disco a lo largo del experimento. Una consulta que vuelca degrada por
una vía que el throttling de CPU y la memoria residente no muestran, y
\texttt{pg\_stat\_database.temp\_bytes} la delata de inmediato.

\paragraph{4. No usar la memoria residente para dimensionar un contenedor de PostgreSQL.}
\texttt{max\_rss\_mb} no cuenta los \texttt{shared\_buffers}, que son memoria
compartida. En este experimento PostgreSQL aparece con 16~MiB residentes y un
conjunto de trabajo de 891~MiB en la misma celda. Un límite de memoria fijado con
la primera cifra mataría el contenedor de inmediato. La métrica útil es
\texttt{working\_set} (o \texttt{memory.current} menos la caché recuperable).

\paragraph{5. Con límites de CPU bajos, el calentamiento tiene que ser largo.}
Con 50 clientes y la caché limpia, la primera consulta de cada cliente tarda de
69 a 110~s en los perfiles B y C. Un calentamiento de 60~s deja ese costo dentro
de la ventana de medición y contamina los percentiles de cola: medido, hasta el
91~\% de las consultas de una ventana eran primeras consultas. Quien mida latencia
de cola bajo contención debe comprobar que su calentamiento cubre la primera
consulta, no suponerlo.

\paragraph{6. Un límite de memoria sin \emph{swap} convierte la saturación en caída, no en lentitud.}
Con \texttt{memory.swap.max = 0}, PostgreSQL con \texttt{Q1} y 50 clientes no se
degradó: fue \textbf{matado por el kernel}, y en tres de las cinco veces antes de
empezar a medir. Para un servicio en producción eso es la diferencia entre
responder despacio y no responder. Si el objetivo es degradar en vez de caer, hace
falta limitar la concurrencia en la capa de conexión antes de que el motor agote su
límite.

\paragraph{7. Medir en un portátil obliga a declarar la frecuencia del procesador.}
Durante la matriz el rendimiento del procesador osciló entre el 86~\% y el
60.8~\% de la frecuencia nominal, y la mayor dispersión observada coincide con una
de esas caídas. En un host dedicado esto no aparecería; en un entorno como este,
hay que medirlo y declararlo, no suponer que la CPU es constante.
